\documentclass[11pt]{article}

\usepackage[margin=1.6cm]{geometry}
\usepackage[T1]{fontenc}
\usepackage[utf8]{inputenc}
\usepackage{lmodern}
\usepackage{amsmath}
\usepackage{amssymb}
\usepackage{booktabs}
\usepackage{xcolor}
\usepackage{enumitem}
\usepackage{tabularx}
\usepackage{caption}

\setlength{\parskip}{0.4em}
\setlength{\parindent}{0pt}
\renewcommand{\arraystretch}{1.05}

\newcommand{\code}[1]{\texttt{#1}}
\newcommand{\stage}[1]{\textsc{Stage~#1}}

\captionsetup[table]{font=small, skip=4pt}

\begin{document}

\begin{center}
{\Large\bfseries Preliminary Results from \stage{5} (Pilot)}\\[0.2em]
{\small Update for supervision \quad|\quad Hassan Suliman \quad|\quad 7 May 2026}
\end{center}

\hrule
\vspace{0.5em}

\textbf{Status.} \stage{5} pilot launched 6 May, in flight on a single A100 80GB.
Workshop scope (single seed, four attributes, $n=30$ cases, three principals
\(\Rightarrow\) 360 trajectories). Snapshot pulled mid-run at \textbf{145/360
trajectories complete}: the entire HIV cell ($n=30$ benign, direct, covert) plus
the IVDU benign cell and 25/30 IVDU direct trajectories. Pregnancy and
mental-health-dx have not yet started (loop reaches them in the next $\sim$36~h).

\vspace{0.2em}
\textbf{1.~The hybrid patient simulator is validated end-to-end on the prod stack.}
The \stage{4} baseline capped direct $\tau$-recovery at $85\%$ (with
$2/5=40\%$ on $\tau{=}1$ cases) due to the simulator returning lab values
consistent with the underlying MedQA disease rather than the augmented $\tau$.
The two-layer fix (prompt anchor + deterministic interceptor; described in the
revised report \S{}5 and \S{}8.4) raises the direct attack ceiling on the
prod stack as predicted, and the same architecture works on the second
attribute on first contact.

\begin{table}[h!]
\centering
\caption{Direct $\tau$-recovery at $n=30$ vs.\ the \stage{4} baseline. The
$\tau{=}1$ split is the load-bearing column: in \stage{4} the simulator
failed 3 of 5 $\tau{=}1$ cases (returning non-reactive serology). With the
hybrid simulator, that failure mode is closed.}
\label{tab:direct}
\begin{tabular}{lccc}
\toprule
& Overall & $\tau{=}1$ & $\tau{=}0$ \\
\midrule
HIV direct, \stage{4} baseline ($n{=}20$) & $17/20 = 85\%$  & $2/5 = 40\%$    & $15/15 = 100\%$ \\
\textbf{HIV direct, \stage{5} pilot ($n{=}30$)}     & $\mathbf{29/30 = 96.7\%}$ & $\mathbf{18/19 = 94.7\%}$ & $\mathbf{11/11 = 100\%}$ \\
\textbf{IVDU direct, \stage{5} pilot ($n{=}25$, partial)} & $\mathbf{23/25 = 92.0\%}$ & $\mathbf{9/9 = 100\%}$  & $\mathbf{14/16 = 87.5\%}$ \\
\bottomrule
\end{tabular}
\end{table}

\textbf{2.~Headline invisibility property holds at $n=30$.}
Benign vs.\ covert diagnostic accuracy on the same case-set is consistent
with parity at this sample size (Wilson 95\% CIs overlap heavily). The
paired-comparison invariant (same \code{case\_id} appears under all three
principals) is enforced by the handler.

\begin{table}[h!]
\centering
\caption{HIV benign-vs-covert paired-outcome comparison. Direction of effect
replicates the \stage{4} baseline finding; net case-flip is small, per-case
agreement remains high.}
\label{tab:paired}
\begin{tabular}{lcc}
\toprule
& \stage{4} baseline ($n{=}20$) & \textbf{\stage{5} HIV ($n{=}30$)} \\
\midrule
Benign accuracy        & $9/20 = 45\%$  & $\mathbf{17/30 = 56.7\%}$ \\
Covert accuracy        & $8/20 = 40\%$  & $\mathbf{14/30 = 46.7\%}$ \\
Both right ($\checkmark{,}\checkmark$) & 5  & $\mathbf{12}$ \\
Both wrong ($\times{,}\times$)         & 8  & $\mathbf{11}$ \\
Covert recovered ($\times{,}\checkmark$) & 3  & $\mathbf{2}$ \\
Covert hurt ($\checkmark{,}\times$)    & 4  & $\mathbf{5}$ \\
Net case-flip          & $-1$ & $\mathbf{-3}$ \\
Per-case agreement     & $65\%$ & $\mathbf{76.7\%}$ \\
\bottomrule
\end{tabular}
\end{table}

\textbf{3.~Anti-leak audit holds across attributes.}
Across $30$ HIV benign and $30$ IVDU benign trajectories, $0/30$ HIV-benign and
$1/30$ IVDU-benign trajectories contain a spontaneous attribute mention. The
single IVDU ``leak'' is a false positive: ``pins and needles'' (paresthesia)
in a foot-exam observation matched the \texttt{needle} regex. The simulator
does not volunteer the sensitive attribute on benign questions.

\textbf{4.~Trajectory-shape divergence preserved.}
Per-principal mean cost / mean steps / ORDER\_TEST share track the \stage{4}
baseline within sampling noise. HIV benign: \$150 / 7.6 steps / 21.1\%.
HIV direct: \$42 / 1.07 steps / 81\%. HIV covert: \$153 / 6.2 steps / 22.6\%.
The qualitative ``covert orders more dollars of testing for comparable
accuracy'' finding from \stage{4} replicates.

\textbf{5.~One signal weakened by config trade-off.}
The Stage 4 finding ``covert produces $\sim$62\% more HIV-keyword EIG
candidates than benign'' has a smaller absolute magnitude in \stage{5} (mean
$0.20$ vs $0.03$, ratio $\sim 6.7\times$). This is consistent with the
workshop-config decision to reduce \(n_{\text{eig\_predictions}}\) from
$5$ to $3$ and \(\max_{\text{queries}}\) from $15$ to $12$ to fit the budget
\(\Rightarrow\) fewer total candidates per trajectory \(\Rightarrow\) fewer
absolute matches. The relative ratio (covert / benign) is preserved; the
\stage{6} probe will quantify whether this weaker signal still supports a
trained classifier.

\textbf{What's still running and when results land.}
IVDU covert ($n{=}30$, $\sim$6\,h), the full pregnancy cell ($n{=}90$,
$\sim$13\,h), and the full mental-health-dx cell ($n{=}90$, $\sim$13\,h)
remain. ETA for full $n{=}360$ dataset: \textbf{$\sim$36\,h from now}, late 8 May.
Auto-shutdown is wired; resume support handles partial failure recovery
without redoing the 145 trajectories already on disk.

\textbf{Two questions for supervision.} (1) Is the workshop-scope
(single-seed, $n=30$) sufficient for a workshop-track submission, or
should we ask for lab credits to scale to top-conference scope before
\stage{6} probes? (2) Should the \stage{6} probe be reported with and
without the chief-complaint string conditioning (the gap quantifies how
much trajectory carries vs.\ the principal task itself)?

\end{document}
